\chapter{Mutual authentication}
\label{ch:auth}

\chapterbrief
  {Prove to the card you hold its key, make it prove the same to you, and establish a session --- without either side ever transmitting the key.}
  {Injecting randomness as a seam.}
  {A real card cryptographically proving it holds your key.}

\section{The problem}

You and the card both --- allegedly --- hold the same sixteen-byte AES key. How
do you each verify that about the other, over a radio link anyone can listen to,
without sending the key?

You cannot send it. You cannot send a hash of it, because an eavesdropper simply
replays the hash. You need each side to demonstrate \emph{use} of the key on
something new every time.

\section{The idea: challenge--response, in both directions}

I send you a random number and ask you to encrypt it. Only someone with the key
can produce the right answer. Because the number is new each time, a recorded
answer is worthless later.

Do it in both directions and you have \textbf{mutual} authentication.

\begin{analogy}{Two spies meeting}
One says ``the nightingale sings at midnight'' --- a line only someone briefed
would recognise.

If only \emph{one} side gives a passphrase, an impostor can collect it and use it
elsewhere. Both sides challenging each other means neither can be impersonated.

The rotation below is the detail that makes it work: you do not send back
\emph{the number you were given}, you send back a \textbf{transformation} of it,
proving you actually decrypted and processed it rather than echoing bytes.
\end{analogy}

\section{The dance}

\begin{center}
\begin{tikzpicture}[font=\scriptsize\sffamily, x=1mm, y=1mm]
  \def\L{0}   \def\R{74}
  % lifelines
  \node[font=\footnotesize\sffamily\bfseries, color=brewernavy] at (\L,4) {READER};
  \node[font=\footnotesize\sffamily\bfseries, color=breweramber!80!black] at (\R,4) {CARD};
  \draw[brewergrey!50] (\L,0) -- (\L,-108);
  \draw[brewergrey!50] (\R,0) -- (\R,-108);

  % msg 1
  \draw[-{Stealth[length=4pt]}, brewernavy] (\L,-6) -- node[above, font=\tiny] {\code{90 71 00 00 02 [keyNo] 00 00}} (\R,-6);
  \node[right=1mm of {(\R,-14)}, anchor=west, font=\tiny\itshape, color=breweramber!80!black, align=left]
       {generates RndB};
  % msg 2
  \draw[-{Stealth[length=4pt]}, breweramber!80!black] (\R,-22) -- node[above, font=\tiny] {\code{E(K, RndB)} \;+\; \hex{91 AF}} (\L,-22);

  \node[left=1mm of {(\L,-32)}, anchor=east, font=\tiny\itshape, color=brewernavy, align=right]
       {decrypt $\rightarrow$ RndB\\ rotate left 1 $\rightarrow$ RndB$'$\\ generate RndA\\ encrypt (RndA $\|$ RndB$'$)};

  % msg 3
  \draw[-{Stealth[length=4pt]}, brewernavy] (\L,-48) -- node[above, font=\tiny] {\code{90 AF 00 00 20 [E(K, RndA||RndB')] 00}} (\R,-48);
  \node[right=1mm of {(\R,-58)}, anchor=west, font=\tiny\itshape, color=breweramber!80!black, align=left]
       {checks RndB$'$ matches\\ $\Rightarrow$ \textbf{reader is authentic}};
  % msg 4
  \draw[-{Stealth[length=4pt]}, breweramber!80!black] (\R,-70) -- node[above, font=\tiny] {\code{E(K, TI || RndA' || caps)} \;+\; \hex{91 00}} (\L,-70);

  \node[left=1mm of {(\L,-84)}, anchor=east, font=\tiny\itshape, color=brewernavy, align=right]
       {check RndA$'$ $=$ rotL(RndA)\\ $\Rightarrow$ \textbf{card is authentic}\\ keep TI; derive $K_{enc}$, $K_{mac}$};

  \draw[brewergreen, dashed] (\L-2,-98) -- (\R+2,-98);
  \node[font=\tiny\bfseries, color=brewergreen] at (37,-102) {SESSION ESTABLISHED};
\end{tikzpicture}
\end{center}

\begin{trap}{Every AES operation here uses an all-zero IV}
Every single one. That is unusual --- normally a fixed IV is a flaw --- and it is
safe \emph{only} because every plaintext contains fresh randomness.

Do not carry the habit anywhere else. In Chapter~\ref{ch:secmsg} the IVs are
derived per command, precisely because the plaintexts there are predictable.
\end{trap}

\paragraph{Rotate left by one byte} means \code{b[1], b[2], ... b[15], b[0]} ---
a \emph{byte} rotation, not a bit rotation:

\begin{trace}
rotateLeft(b0b1b2b3b4b5b6b7b8b9babbbcbdbebf)
         = b1b2b3b4b5b6b7b8b9babbbcbdbebfb0
\end{trace}

\paragraph{TI}, the Transaction Identifier, is four bytes the card generates
fresh for this session. From Chapter~\ref{ch:secmsg} on, it goes into every MAC
and every IV, binding all subsequent traffic to \emph{this} authentication. It is
the session's fingerprint.

\subsection{The step-2 reply, byte by byte}

The card returns 32 encrypted bytes; decrypted they are:

\begin{center}\small
\begin{tabular}{@{}lll@{}}
\toprule
\thead{Offset} & \thead{Size} & \thead{Contents} \\
\midrule
\code{[0..4)}   & 4  & TI --- the transaction identifier \\
\code{[4..20)}  & 16 & RndA$'$ --- your RndA, rotated left \\
\code{[20..26)} & 6  & PDcap2 --- card capabilities \\
\code{[26..32)} & 6  & PCDcap2 --- reader capabilities echoed back \\
\bottomrule
\end{tabular}
\end{center}

You need the first two. The capabilities you can ignore.

\section{The testability problem, and the seam that solves it}

Here is a genuine engineering puzzle. This function calls a random number
generator, so it \textbf{produces different output every run}. How do you write a
known-answer test for it?

The answer is the same move as Chapter~\ref{ch:spi}: \textbf{the thing that
varies becomes a parameter.} Randomness is a dependency like any other. Inject
it, and a test can inject a \emph{fixed} ``random'' number, making the whole
handshake deterministic and checkable.

\begin{jcode}
package com.midnightbrewer.hardware;

import java.security.SecureRandom;

/**
 * A source of random bytes -- the seam that makes the auth handshake testable.
 *
 * Production uses secure(). A test injects a fake returning a FIXED RndA, and
 * suddenly the whole handshake is a known-answer test.
 *
 * @FunctionalInterface = exactly one abstract method, so a fake can be written
 * as a one-line lambda:   RandomSource fixed = n -> myFixedBytes;
 */
@FunctionalInterface
public interface RandomSource {

    byte[] nextBytes(int n);

    /** The real one, backed by a cryptographically secure RNG. */
    static RandomSource secure() {
        SecureRandom rng = new SecureRandom();
        return n -> {
            byte[] out = new byte[n];
            rng.nextBytes(out);
            return out;
        };
    }
}
\end{jcode}

\begin{trap}{\texttt{SecureRandom}, never \texttt{java.util.Random}}
\code{Random} is a linear congruential generator: observe a couple of outputs and
you can predict every future one. For a nonce, that is the entire security of the
protocol.

\code{SecureRandom} draws from the operating system's entropy pool. This is a
one-word difference with a catastrophic failure mode, and it is exactly the kind
of thing that ships.
\end{trap}

\begin{concept}{A seam is not a licence}
The \code{RandomSource} seam exists so \emph{tests} can pin randomness. If
production code ever passes anything but \code{RandomSource.secure()}, the
protocol is broken.

Test doubles live in \code{src/test} and that is not an accident: the build
system enforces that a fixed-nonce implementation cannot reach the kiosk.
\end{concept}

\section{Build it}

\begin{jcode}
package com.midnightbrewer.hardware;

import com.midnightbrewer.hardware.crypto.Aes;
import com.midnightbrewer.hardware.crypto.DesfireSession;
import java.util.Arrays;

/**
 * A DESFire card, spoken to over an ApduChannel.
 *
 * COMPOSITION: it HAS-A channel and HAS-A random source, both injected. It
 * knows nothing about SPI, the RC522, or ISO-DEP framing -- only APDUs in and
 * responses out. That is why it can be tested with no hardware at all.
 */
public class DesfireCard {

    private final ApduChannel channel;
    private final RandomSource randomSource;

    private byte[] ti;                 // transaction identifier, set at auth
    private DesfireSession session;    // the derived keys, set at auth

    public DesfireCard(ApduChannel channel, RandomSource randomSource) {
        this.channel = channel;
        this.randomSource = randomSource;
    }

    /**
     * Wrap a DESFire command in an ISO-7816 APDU and send it.
     *   90 [cmd] 00 00 [Lc] [data...] 00
     */
    private byte[] wrapApdu(byte cmd, byte[] data) throws SpiException {
        byte[] apdu = new byte[data.length + 6];
        apdu[0] = (byte) 0x90;              // CLA: DESFire native, wrapped
        apdu[1] = cmd;                      // INS
        apdu[2] = 0x00;                     // P1
        apdu[3] = 0x00;                     // P2
        apdu[4] = (byte) data.length;       // Lc
        System.arraycopy(data, 0, apdu, 5, data.length);
        apdu[apdu.length - 1] = 0x00;       // Le
        return channel.transceive(apdu);
    }

    /** Rotate left by one BYTE: b[1..n-1] followed by b[0]. */
    private static byte[] rotateLeft(byte[] array) {
        // TODO
        return array;
    }

    /** SelectApplication (0x5A). The AID goes on the wire LSB-first. */
    public void selectApplication(byte[] aid) throws SpiException {
        if (aid.length != 3) {
            throw new IllegalArgumentException("AID must be 3 bytes");
        }
        byte[] response = wrapApdu((byte) 0x5A, aid);
        if (response.length < 2) {
            throw new SpiException("SelectApplication failed");
        }
    }

    // =================================================================
    // AuthenticateEV2First (0x71). All AES here uses an ALL-ZERO IV.
    // =================================================================
    public void authenticateEv2First(byte keyNo, byte[] key) throws SpiException {
        if (key.length != 16) {
            throw new IllegalArgumentException("key must be 16 bytes");
        }

        // ---- step 1: ask for the challenge ----
        // TODO 1: byte[] r1 = wrapApdu((byte) 0x71, new byte[]{ keyNo, 0x00 });
        //
        //   TWO data bytes: the key number AND a 0x00 "length of capabilities"
        //   byte. Sending only keyNo gets you 91 7E (length error) and it is a
        //   genuinely annoying hour to find.
        //
        // TODO 2: strip the 2 status bytes -> encRndB (16 bytes)
        // TODO 3: byte[] rndB = Aes.cbcDecrypt(key, new byte[16], encRndB);

        // ---- step 2: answer, and challenge back ----
        // TODO 4: byte[] rndBRot = rotateLeft(rndB);
        //         byte[] rndA    = randomSource.nextBytes(16);
        // TODO 5: build a 32-byte buffer: rndA then rndBRot.
        //         byte[] encAB = Aes.cbcEncrypt(key, new byte[16], that);
        //
        //   ONE 32-byte CBC run, not two 16-byte ones. In CBC the second block
        //   chains off the first's ciphertext -- encrypting them separately
        //   gives a different, wrong answer.
        //
        // TODO 6: byte[] r2 = wrapApdu((byte) 0xAF, encAB);
        //         strip the status bytes -> encResp (32 bytes)

        // ---- step 3: verify the card, keep the session ----
        // TODO 7: byte[] plain = Aes.cbcDecrypt(key, new byte[16], encResp);
        // TODO 8: byte[] rndARot = Arrays.copyOfRange(plain, 4, 20);
        //         if (!Arrays.equals(rndARot, rotateLeft(rndA)))
        //             throw new SpiException("authentication failed: RndA mismatch");
        //
        //   THIS CHECK IS THE WHOLE POINT. Skip it and you have not
        //   authenticated the card -- you have merely talked to something.
        //
        // TODO 9: this.ti = Arrays.copyOfRange(plain, 0, 4);
        //         this.session = new DesfireSession(key, rndA, rndB);
    }

    public byte[] getTi() { return ti; }
}
\end{jcode}

\section{Verify --- the full handshake, offline}

The elegant part: your test class can \textbf{be} the fake card and the fake
random source at once.

\begin{jcode}
package com.midnightbrewer.hardware;

import java.util.Arrays;
import java.util.HexFormat;
import org.junit.jupiter.api.Test;
import static org.junit.jupiter.api.Assertions.assertArrayEquals;

/**
 * The test IS its own fake card + fake RNG: it implements both seams, so
 * new DesfireCard(this, this) plugs the test straight into the card.
 */
public class DesfireCardTest implements RandomSource, ApduChannel {

    byte[] key   = HexFormat.of().parseHex("00112233445566778899aabbccddeeff");
    byte[] rndA  = HexFormat.of().parseHex("a0a1a2a3a4a5a6a7a8a9aaabacadaeaf");
    byte[] encAB = HexFormat.of().parseHex(
            "cf086a82c0b745a749daabb28a7a8db3"
          + "32272f9e5fdc37a4746064f25b7c2ca7");
    byte[] tiExpected = HexFormat.of().parseHex("11223344");

    // Replies MUST carry the trailing status word -- the card code strips the
    // last two bytes off every reply. Without them you chop real ciphertext
    // and decrypt garbage.
    byte[] selectReply = HexFormat.of().parseHex("9100");
    byte[] step1Reply  = HexFormat.of().parseHex(
            "eab6822fe368d5bc9895fb2558b38dde" + "91af");
    byte[] step2Reply  = HexFormat.of().parseHex(
            "4e08e447d74810052caa1c5b6c60343f"
          + "51f6e901ad171e677b100d3a34c925e4" + "9100");

    @Override
    public byte[] nextBytes(int n) {
        return rndA;                        // pin RndA -> deterministic
    }

    @Override
    public byte[] transceive(byte[] apdu) throws SpiException {
        switch (apdu[1]) {                  // dispatch on INS
            case (byte) 0x5A:
                return selectReply;
            case (byte) 0x71:
                return step1Reply;
            case (byte) 0xAF:
                // assert on what we SENT: encAB sits at [5 .. len-1)
                assertArrayEquals(encAB,
                        Arrays.copyOfRange(apdu, 5, apdu.length - 1),
                        "step-2 ciphertext is wrong");
                return step2Reply;
            default:
                throw new SpiException("unexpected APDU");
        }
    }

    @Test
    public void authenticateEv2First_matchesReplayVectors() throws SpiException {
        DesfireCard card = new DesfireCard(this, this);

        card.selectApplication(new byte[]{ 0x01, 0x02, 0x03 });
        card.authenticateEv2First((byte) 0x00, key);  // no throw == RndA' verified

        assertArrayEquals(tiExpected, card.getTi());
    }
}
\end{jcode}

\begin{concept}{Why this test is strong}
It checks \emph{both} directions:

\begin{itemize}
  \item the \hex{0xAF} case asserts the exact ciphertext \textbf{you sent} ---
        proving your rotate, your concatenation and your CBC encryption are all
        correct;
  \item reaching the end without throwing proves your RndA$'$ verification passed;
  \item the TI assertion proves you sliced the reply correctly.
\end{itemize}

A test that only checked the return value would pass even if you had skipped the
RndA$'$ comparison entirely.
\end{concept}

\begin{verifybox}
\begin{shell}
mvn test -Dtest=DesfireCardTest
\end{shell}
The full vector set is in Appendix~\ref{app:vectors}.
\end{verifybox}

\section{Then: on the real card}

\begin{jcode}
DesfireCard desfire = new DesfireCard(channel, RandomSource.secure());
desfire.selectApplication(new byte[]{ 0x03, 0x02, 0x01 });  // AID 0x010203, LSB-first
byte[] appKey = new byte[16];
Arrays.fill(appKey, (byte) 0x22);                           // your app's user key
desfire.authenticateEv2First((byte) 0x02, appKey);          // key slot 2
System.out.println("MUTUAL AUTH OK. TI = " + hex(desfire.getTi()));
\end{jcode}

\begin{verifybox}
\begin{trace}
MUTUAL AUTH OK. TI = 8071F582
\end{trace}
Run it twice. \textbf{The TI will be different every time} --- it is freshly
generated per session. Seeing it change is your proof that this was a live
authentication and not a replay of your test vectors.
\end{verifybox}

\begin{trap}{A factory-blank card will fail here}
Its default master key is DES, not AES, and there is no application yet. The card
must be personalised first --- that is Chapter~\ref{ch:android}'s job.

If you do not have a provisioned card, use the top-up app's setup screen before
attempting this chapter on hardware.
\end{trap}

\section{Traps}

\begin{traptable}
\hex{91 7E} (length error) on step 1 & Sent only \code{keyNo}. It must be \code{\{ keyNo, 0x00 \}}. \\
\hex{91 AE} (auth error) on step 2 & Wrong key, wrong key number, or wrong application selected. \\
RndA mismatch but vectors pass & Check you rotated \emph{left}, and by a byte not a bit. \\
Test decrypts garbage & Fake replies missing their \hex{91 xx} status word. \\
Ciphertext differs from the vector & You encrypted RndA and RndB$'$ as two separate CBC calls. \\
Auth ``succeeds'' always & You skipped the RndA$'$ comparison. Verify you actually throw. \\
\end{traptable}

\section{Commit}

\begin{shell}
git add -A
git commit -m "M7c: AuthenticateEV2First, verified offline against replay vectors"
\end{shell}
